\documentclass[sigconf]{acmart}

\input{preamble}

\begin{document}

\title{Augmenting LLM Code Translation with Compiler Analysis for C to Triton Kernel Generation}

\author{Xiao Qin}
\affiliation{
  \institution{University of Leeds}
  \city{Leeds}
  \country{United Kingdom}
}
\email{ppnh4382@leeds.ac.uk}

\author{Chunwei Xia}
\affiliation{
  \institution{University of Leeds}
  \city{Leeds}
  \country{United Kingdom}
}

\author{Zheng Wang}
\affiliation{
  \institution{University of Leeds}
  \city{Leeds}
  \country{United Kingdom}
}




\begin{abstract}
Emerging programming models like Triton enable developers to better exploit modern accelerators, but translating legacy code to Triton remains challenging. While Large Language Models (LLMs) show promise for code translation, they often generate incorrect or suboptimal implementations due to a lack of precise parallelization reasoning.
%
We present \sysname, a compiler-assisted LLM framework for translating C loop nests to Triton kernels. \sysname extracts parallelization properties using compiler-based analysis and encodes them as structured analysis facts in the LLM prompt, enabling correct and performance-aware parallel code synthesis. A profiling feedback loop further improves performance by feeding GPU hardware utilization metrics back to the LLM for iterative optimization.
%
On PolyBench/C, \sysname achieves 97\% correctness, outperforming both an autonomous LLM agent with tool use (80\%) and an unguided LLM baseline (93\%). At 8$\times$ problem sizes, the generated Triton kernels achieve a median 24.6$\times$ speedup, compared to 20.7$\times$ for the agent and 10.8$\times$ for the unguided baseline. Profiling feedback improves 9 of 29 kernels at standard sizes and 12 of 21 at larger sizes. \sysname also passes 95\% of 151 TSVC kernels and generalizes to 6 of 8 Rodinia/ECP application kernels. These results show that compiler analysis outperforms autonomous LLM agents for kernel generation correctness and performance.


% Emerging programming models like Trition allow programs to take advantage of parallel computing hardware,  but translating legacy code to Trition remains a challenge. Large language models (LLMs) can transpile between programming langguages,  but without
% understanding parallelization constraints they frequently produce
% incorrect or inefficient implementations.  We present \sysname, a
% system that extracts parallelization properties from C loop nests
% using polyhedral analysis (\pet) and LLVM's dependence analysis, then
% encodes them as structured natural-language guidance in the LLM prompt.
% On 30 Polybench/C kernels, \sysname achieves 100\% correctness
% (vs.\ 93\% without analysis) with 2.15$\times$ median speedup over
% single-threaded C, rising to 30.73$\times$ at 8$\times$ problem sizes.
% On 151 TSVC loop kernels, \sysname passes 150/151 (99\%);
% at 32\,MB per array, the generated GPU kernels achieve 46$\times$
% median speedup over single-threaded C and 9$\times$ over
% 80-thread OpenMP.
% The approach generalizes to 8 application kernels from Rodinia and
% ECP proxy applications, achieving up to 25.8$\times$ speedup.
\end{abstract}

\begin{CCSXML}
<ccs2012>
  <concept>
    <concept_id>10011007.10011006.10011041</concept_id>
    <concept_desc>Software and its engineering~Compilers</concept_desc>
    <concept_significance>500</concept_significance>
  </concept>
  <concept>
    <concept_id>10010147.10010169.10010170</concept_id>
    <concept_desc>Computing methodologies~Parallel computing methodologies</concept_desc>
    <concept_significance>300</concept_significance>
  </concept>
</ccs2012>
\end{CCSXML}

\ccsdesc[500]{Software and its engineering~Compilers}
\ccsdesc[300]{Computing methodologies~Parallel computing methodologies}

\keywords{LLM code generation, GPU kernels, Triton, polyhedral analysis}

\maketitle

%% ===================================================================
\section{Introduction}
\label{sec:intro}

Accelerators such as GPUs are now essential for high-performance scientific computing, yet much existing software is still written in legacy C code that is not designed for parallel execution on modern accelerators. Translating these legacy programs to accelerator-friendly implementations is challenging, as it requires identifying parallelism, restructuring computation, and managing synchronization, tasks that require expertise in programming models and hardware.


Emerging programming models such as Triton~\cite{tillet2019triton} simplify GPU programming by providing higher-level abstractions for writing efficient kernels. However, mapping legacy C code to Triton still requires careful reasoning about loop-level parallelism, data dependencies, and execution order. Developers must determine which loop dimensions can be safely parallelized, how memory accesses should be organized, and when synchronization is required - decisions that are subtle and error-prone in practice.


Large Language Models (LLMs) have recently shown promise in code translation, including generating GPU kernels from high-level descriptions~\cite{chen2024chatgpt, ouyang2024kernelbench}. However, LLMs lack precise reasoning about program semantics, especially data dependencies. As a result, they often produce incorrect or inefficient code, including unsafe parallelization that leads to data races, suboptimal kernel structures, and missing synchronization.


In this paper, we present \sysname, a system that augments LLM-based translation with compiler analysis. \sysname extracts parallelization properties from C loop nests using polyhedral analysis (\pet)~\cite{verdoolaege2012polyhedral} and LLVM dependence analysis~\cite{llvm}, and encodes them as structured guidance in the LLM prompt. This allows the model to generate correct and performance-aware Triton code without requiring it to perform dependence analysis itself. 

We evaluate \sysname on PolyBench/C, TSVC, and application kernels on an NVIDIA RTX~3090 GPU. \sysname achieves 97\% correctness on PolyBench/C and 95\% on TSVC, improving over unguided baselines. At 8$\times$ problem sizes, the generated kernels achieve a median 24.6$\times$ speedup. We further demonstrate generalization on 8 Rodinia and ECP application kernels. This paper makes the following contributions:

\begin{itemize}
\item A compiler-assisted LLM framework for translating C loop nests to Triton, with a unified analysis module that derives all parallelization constraints from a single polyhedral analysis invocation and encodes them as structured facts.
\item A profiling-guided optimization loop that feeds GPU hardware utilization metrics (from NVIDIA Nsight Compute) back to the LLM, enabling iterative performance improvement while preserving correctness.
\item Empirical evaluation on 189 kernels across PolyBench/C, TSVC, Rodinia, and ECP benchmarks, showing that compiler analysis improves correctness and profiling feedback significantly improves performance.
\end{itemize}


\section{Our Approach}
\label{sec:method}
\input{workflow}

Figure~\ref{fig:pipeline} shows the \sysname workflow.
Given a C source file, \sysname extracts loop nests, analyzes their parallelization properties, constructs a structured prompt, invokes an LLM to generate Triton code, and optionally refines performance using GPU profiling feedback.

\subsection{Analysis and Prompt Construction}
\sysname calls PET~\cite{verdoolaege2012polyhedral} once per kernel to extract a polyhedral representation (iteration domains, access relations, schedules). From this single invocation, it derives all parallelization constraints in a unified JSON schema:
\emph{(i)}~for each loop dimension, whether it can be safely parallelized (by checking for cross-iteration dependences via ISL relation composition);
\emph{(ii)}~WAR (write-after-read) dependences that require array cloning or sequential execution;
\emph{(iii)}~reduction patterns (scalar accumulation + array read) mapped to Triton primitives;
\emph{(iv)}~memory access patterns (stencil, triangular, cross-phase) that affect kernel structure.
When polyhedral analysis fails (e.g., non-affine accesses), \sysname falls back to LLVM's \texttt{DependenceAnalysis}.

A single, pattern-agnostic function renders the analysis into structured text: it emits whatever constraints were found without pattern-specific templates. The prompt includes the C source, analysis facts (e.g., ``\texttt{j} is safe; \texttt{t} has cross-phase deps''), and Triton coding rules. The LLM receives facts and decides the implementation strategy.

\subsection{Generation and Iterative Refinement}
\sysname invokes Claude Sonnet~4~\cite{anthropic2025claude} and validates the generated kernel against the C reference using identical inputs and FP32 tolerance checks. On failure, errors are classified (compilation, numerical, missing barriers, low parallelism) and a targeted retry prompt is issued, up to 10 attempts with a context reset at attempt~6.

\subsection{Profiling-Guided Optimization}
\label{sec:profiling}
After correctness validation, \sysname optionally profiles the kernel with NVIDIA Nsight Compute (NCU), collecting SM throughput, memory throughput, occupancy, and cache hit rates. It classifies the bottleneck (\emph{compute-bound}, \emph{memory-bound}, or \emph{latency-bound}) and feeds the metrics back to the LLM with bottleneck-specific optimization guidance. The optimized kernel is re-validated for correctness and kept only if faster. This loop runs up to 3 iterations, with NCU results cached when the implementation is unchanged.


\section{Evaluation Setup}

\cparagraph{Platforms.}
All experiments are conducted on a GPU workstation equipped with an NVIDIA RTX~3090 GPU (24\,GB memory, 936\,GB/s memory bandwidth) and a dual-socket Intel Xeon Gold 5218R CPU (80 logical cores).
We use PyTorch~2.5.1, Triton~3.1.0, CUDA~12.4, GCC~11.4.0, and Claude Sonnet~4~\cite{anthropic2025claude} as the underlying LLM.

\cparagraph{Benchmarks.}
We evaluate on three benchmark suites:
(1)~\emph{PolyBench/C}~\cite{polybench}: 30 numerical kernels (stencils, linear algebra, data mining) at default sizes ($N$=60--120) and 8$\times$ scale ($N$=480--960);
(2)~\emph{TSVC}~\cite{callahan1988tsvc}: 151 loop kernels covering linear traversals, reductions, indirect addressing, and loop-carried dependences;
(3)~8 application kernels from Rodinia~\cite{che2009rodinia} and ECP proxy applications~\cite{heroux2009mantevo}.

\cparagraph{Configurations.}
On PolyBench/C (our primary benchmark), we compare three configurations:
\begin{itemize}
\item \emph{Agent}: An autonomous LLM agent with tool use (write code, run tests, run benchmarks). The agent sees full error output and decides its own retry strategy. No compiler analysis. Up to 10 correctness turns + 3 optimization turns.
\item \emph{NA} (No Analysis): The LLM receives only the C source code with classified error feedback on retry. Same turn budget as the agent.
\item \emph{WA+Prof} (With Analysis + Profiling): Our full method. The LLM receives compiler-derived analysis facts. After correctness, NCU profiling feedback guides up to 3 optimization iterations.
\end{itemize}
On TSVC and application kernels, we report WA+Prof only (our method).

\cparagraph{Evaluation methodology.}
Speedups are measured over 50 iterations after 5 warmup iterations. Each iteration includes array re-initialization (GPU \texttt{clone()} for Triton, CPU \texttt{numpy.copy()} for the C reference) to ensure in-place kernels start from identical data. Triton timing includes \texttt{cuda.synchronize()}.

%% ===================================================================
\section{Experimental Results}
\label{sec:eval}

\subsection{PolyBench/C (30 Kernels)}

\begin{table}[t]
  \caption{PolyBench/C results at 1$\times$ sizes. Agent = autonomous LLM with tool use (no analysis); NA = LLM with classified error feedback (no analysis); WA+Prof = our method (compiler analysis + NCU profiling feedback).}
  \label{tab:results_1x}
  \centering\small
  \begin{tabular}{lccc}
    \toprule
    \textbf{Metric} & \textbf{Agent} & \textbf{NA} & \textbf{WA+Prof} \\
    \midrule
    Pass rate          & 24/30 (80\%) & 28/30 (93\%) & 29/30 (97\%) \\
    First-try passes   & 11           & 8            & 12           \\
    Median speedup     & 0.60$\times$ & 0.88$\times$ & 0.76$\times$ \\
    Mean speedup       & 0.87$\times$ & 1.64$\times$ & 1.24$\times$ \\
    Kernels $>1\times$ & 7/23         & 12/28        & 11/29        \\
    \bottomrule
  \end{tabular}
\end{table}

\begin{table}[t]
  \caption{PolyBench/C results at 8$\times$ sizes ($N$: 60$\to$480, matrices: 120$\to$960). Speedups vs.\ single-threaded C.}
  \label{tab:results_8x}
  \centering\small
  \begin{tabular}{lccc}
    \toprule
    \textbf{Metric} & \textbf{Agent} & \textbf{NA} & \textbf{WA+Prof} \\
    \midrule
    Pass rate          & 17/30 (57\%) & 26/30 (87\%) & 27/30 (90\%) \\
    Median speedup     & 20.7$\times$ & 10.8$\times$ & 24.6$\times$ \\
    Kernels $>1\times$ & 13/15        & 14/20        & 18/21        \\
    \bottomrule
  \end{tabular}
\end{table}

Tables~\ref{tab:results_1x} and~\ref{tab:results_8x} compare three configurations on PolyBench/C.
Our method (WA+Prof) achieves the highest correctness at both scales: 97\% at 1$\times$ and 90\% at 8$\times$.
The autonomous agent baseline, despite having tool access (write code, run tests, run benchmarks), achieves only 80\% at 1$\times$ and 57\% at 8$\times$---the weakest of all three.
This shows that autonomous tool use without compiler analysis guidance leads to poor retry strategies: the agent wastes turns on ineffective fixes without understanding the root cause of parallelization failures.

At 1$\times$ sizes, median speedups are limited by GPU launch overhead at small $N$.
At 8$\times$ sizes, WA+Prof achieves 24.6$\times$ median, outperforming NA (10.8$\times$) and the agent (20.7$\times$).
Top kernels: \texttt{3mm} (2041$\times$), \texttt{2mm} (2004$\times$), \texttt{gemm} (451$\times$).

\subsection{Profiling-Guided Optimization}

\begin{table}[t]
  \caption{Representative profiling feedback improvements.
    ``Before'' = first correct kernel; ``After'' = best after NCU-guided iterations.}
  \label{tab:profiling}
  \centering\small
  \begin{tabular}{lccl}
    \toprule
    \textbf{Kernel} & \textbf{Before} & \textbf{After} & \textbf{Scale} \\
    \midrule
    3mm         & 0.10$\times$ & 4.01$\times$ & 1$\times$ \\
    covariance  & 0.11$\times$ & 3.66$\times$ & 1$\times$ \\
    2mm         & 1.48$\times$ & 2.97$\times$ & 1$\times$ \\
    syrk        & 1.65$\times$ & 2.07$\times$ & 1$\times$ \\
    \midrule
    gemm        & 120$\times$  & 451$\times$  & 8$\times$ \\
    correlation & 12$\times$   & 51$\times$   & 8$\times$ \\
    atax        & 0.4$\times$  & 15$\times$   & 8$\times$ \\
    syrk        & 24$\times$   & 62$\times$   & 8$\times$ \\
    \bottomrule
  \end{tabular}
\end{table}

Profiling feedback improved 9 of 29 kernels at 1$\times$ and 12 of 21 benchmarked kernels at 8$\times$ (Table~\ref{tab:profiling}).
The gains are largest on latency-bound kernels where the initial implementation has insufficient parallelism: the NCU metrics reveal low SM and memory utilization, and the LLM restructures the kernel accordingly.
Correctness is preserved: each optimized kernel is re-validated, and attempts that break correctness are rejected.

\subsection{TSVC (151 Kernels)}

\sysname passes 143/151 TSVC kernels (95\%), 91 on first attempt,
with a mean speedup of 2.78$\times$ and maximum of 267$\times$ over single-threaded C.

\subsection{Application Kernels}

\begin{table}[t]
  \caption{Application kernels: Triton speedup vs.\ single-threaded C.}
  \label{tab:realworld}
  \centering\small
  \begin{tabular}{llcc}
    \toprule
    \textbf{Source} & \textbf{Kernel} & \textbf{WA/ST} & \textbf{Pass?} \\
    \midrule
    ECP     & LJ force   & 24.1$\times$ & Y \\
    ECP     & SpMV       & 7.1$\times$  & Y \\
    Rodinia & SRAD       & 12.8$\times$ & Y \\
    Rodinia & pathfinder & 3.6$\times$  & Y \\
    Rodinia & lud        & 1.3$\times$  & Y \\
    Rodinia & gaussian   & 0.6$\times$  & Y \\
    \midrule
    Rodinia & hotspot    & ---          & N \\
    Rodinia & lavaMD     & ---          & N \\
    \bottomrule
  \end{tabular}
\end{table}

We apply \sysname to 3 Rodinia and 5 ECP/Rodinia application kernels (Table~\ref{tab:realworld}).
6 of 8 pass correctness.
The highest speedups come from compute-intensive kernels:
\texttt{LJ\_force} (24.1$\times$) and \texttt{SRAD} (12.8$\times$).
Two kernels fail: \texttt{hotspot} (complex 2D stencil with boundary conditions) and \texttt{lavaMD} (irregular particle interactions), both requiring domain-specific optimizations beyond what the unified analysis provides.



%% ===================================================================
\section{Related Work}
\label{sec:related}

\paragraph{LLM Code Generation for GPUs.}
Chen et al.~\cite{chen2024chatgpt} find frequent correctness issues in ChatGPT-generated CUDA kernels.
KernelBench~\cite{ouyang2024kernelbench} reports 20--40\% functional correctness for LLM-generated GPU kernels without guidance.
Recent works improve generation through RL fine-tuning: Dr.~Kernel~\cite{liu2026drkernel} uses profiling-based rewards, Tehrani et al.~\cite{tehrani2026ftgpt5} RL-fine-tune GPT-5, and DICE~\cite{bai2026dice} applies RL to diffusion LLMs.
Search-based methods keep the LLM frozen: AVO~\cite{chen2026avo} uses LLM agents as evolutionary operators, K-Search~\cite{cao2026ksearch} decouples planning from code instantiation, and OptiML~\cite{bhattacharjee2026optiml} uses MCTS with profiler feedback.
KernelBlaster~\cite{dong2026kernelblaster} enables cross-task learning via persistent memory.
For non-GPU targets, AscendKernelGen~\cite{cao2026ascendkernelgen} and AscendCraft~\cite{wen2026ascendcraft} address NPU kernels.
These works focus on \emph{how} to generate/optimize kernels; our work is orthogonal, addressing \emph{what} parallelism is safe via compiler analysis. The two directions are complementary: compiler-derived constraints could serve as correctness rewards in RL frameworks or safety priors in search.

\paragraph{Polyhedral Compilation and Auto-Scheduling.}
PPCG~\cite{verdoolaege2013ppcg} generates CUDA from affine loop nests but produces low-level code that lacks Triton's block-level optimizations.
TVM~\cite{chen2018tvm} and Halide~\cite{ragan2013halide} use auto-scheduling with learned cost models.
\sysname is complementary: it uses polyhedral analysis to extract \emph{what} parallelism is safe, then delegates the \emph{how} to the LLM, which can handle non-affine constructs, irregular access patterns, and Triton abstractions that polyhedral compilers cannot target directly.


%% ===================================================================
\section{Conclusion \& Future Work}
\label{sec:conclusion}
We presented \sysname, a system that augments LLM-based code generation with compiler analysis for translating C loop nests to Triton GPU kernels. By deriving all parallelization constraints from a single polyhedral analysis invocation and rendering them as structured facts, \sysname provides the LLM with precise correctness guidance without pattern-specific prompt engineering.

Our experiments demonstrate that compiler-derived analysis significantly improves both correctness and performance over unguided LLM baselines and autonomous LLM agents. Notably, an autonomous agent with full tool access (code execution, testing, debugging) achieves lower correctness than our analysis-guided approach, showing that structured compiler knowledge is more valuable than autonomous trial-and-error for parallel code generation. The profiling feedback loop further improves performance by identifying and addressing hardware utilization bottlenecks through iterative optimization.

Several directions remain for future work.
First, extending the agent baseline comparison to TSVC and application benchmarks would provide a more comprehensive evaluation of when compiler analysis outperforms autonomous tool use.
Second, our current correctness validation relies on input/output comparison, which cannot guarantee equivalence for all inputs. Integrating formal equivalence checking for affine programs~\cite{verdoolaege2012equiv,barthou2012widening,lefevre2023verified} would provide stronger correctness guarantees, particularly for kernels with complex recurrence patterns where test-based validation is insufficient.
Third, following the decoupled planning-and-implementation strategy shown effective in recent kernel search methods~\cite{cao2026ksearch}, we plan to separate high-level parallelization strategy selection (guided by compiler analysis) from low-level Triton code instantiation, enabling the system to explore non-monotonic optimization paths without abandoning promising strategies due to buggy intermediate implementations.
Fourth, analyzing the \emph{generated} Triton code using Triton's own compiler IR (e.g., memory coalescing analysis, occupancy estimation) and feeding this back to the LLM would create a closed-loop optimization cycle that leverages target-language-specific information unavailable from the source-level analysis alone.

%% ===================================================================
%% References
\bibliographystyle{ACM-Reference-Format}
\begin{thebibliography}{99}

\bibitem{tillet2019triton}
P.~Tillet, H.~T. Kung, and D.~Cox.
\newblock Triton: An intermediate language and compiler for tiled neural
  network computations.
\newblock In \emph{Proc. MAPL}, 2019.

\bibitem{chen2024chatgpt}
Y.~Chen, S.~Huang, D.~Zheng, et~al.
\newblock {ChatGPT} for {CUDA} programming: Evaluating {LLM}-generated {GPU}
  code.
\newblock \emph{arXiv preprint arXiv:2404.12000}, 2024.

\bibitem{ouyang2024kernelbench}
A.~Ouyang, S.~Zheng, et~al.
\newblock {KernelBench}: Can {LLMs} write efficient {GPU} kernels?
\newblock \emph{arXiv preprint arXiv:2404.12489}, 2024.

\bibitem{verdoolaege2012polyhedral}
S.~Verdoolaege and T.~Grosser.
\newblock Polyhedral extraction tool.
\newblock In \emph{Proc. IMPACT}, 2012.

\bibitem{llvm}
C.~Lattner and V.~Adve.
\newblock {LLVM}: A compilation framework for lifelong program analysis \&
  transformation.
\newblock In \emph{Proc. CGO}, 2004.

\bibitem{polybench}
L.-N. Pouchet.
\newblock Polybench/C: The polyhedral benchmark suite, version 4.2.1.
\newblock \url{https://web.cs.ucla.edu/~pouchet/software/polybench/}, 2016.

\bibitem{anthropic2025claude}
Anthropic.
\newblock Claude {Sonnet} 4 model card.
\newblock \url{https://www.anthropic.com/claude}, 2025.

\bibitem{verdoolaege2010isl}
S.~Verdoolaege.
\newblock isl: An integer set library for the polyhedral model.
\newblock In \emph{Proc. ICMS}, 2010.

\bibitem{chen2018tvm}
T.~Chen, T.~Moreau, Z.~Jiang, et~al.
\newblock {TVM}: An automated end-to-end optimizing compiler for deep
  learning.
\newblock In \emph{Proc. OSDI}, 2018.

\bibitem{ragan2013halide}
J.~Ragan-Kelley, C.~Barnes, A.~Adams, et~al.
\newblock Halide: A language and compiler for optimizing parallelism,
  locality, and recomputation in image processing pipelines.
\newblock In \emph{Proc. PLDI}, 2013.

\bibitem{bondhugula2008pluto}
U.~Bondhugula, A.~Hartono, J.~Ramanujam, and P.~Sadayappan.
\newblock A practical automatic polyhedral parallelizer and locality
  optimizer.
\newblock In \emph{Proc. PLDI}, 2008.

\bibitem{verdoolaege2013ppcg}
S.~Verdoolaege, J.~Carlos~Juega, A.~Cohen, et~al.
\newblock Polyhedral parallel code generation for {CUDA}.
\newblock \emph{ACM TACO}, 9(4), 2013.

\bibitem{che2009rodinia}
S.~Che, M.~Boyer, J.~Meng, et~al.
\newblock Rodinia: A benchmark suite for heterogeneous computing.
\newblock In \emph{Proc. IISWC}, 2009.

\bibitem{heroux2009mantevo}
M.~A.~Heroux, D.~W.~Doerfler, P.~S.~Crozier, et~al.
\newblock Improving performance via mini-applications.
\newblock Sandia National Laboratories Technical Report SAND2009-5574, 2009.

\bibitem{callahan1988tsvc}
D.~Callahan, J.~Dongarra, and D.~Levine.
\newblock Vectorizing compilers: A test suite and results.
\newblock In \emph{Proc. Supercomputing}, 1988.

\bibitem{liu2026drkernel}
W.~Liu, J.~Xu, Y.~Li, et~al.
\newblock Dr.~Kernel: Reinforcement learning done right for {Triton} kernel generations.
\newblock \emph{arXiv preprint arXiv:2602.05885}, 2026.

\bibitem{tehrani2026ftgpt5}
A.~Tehrani, Y.~Emara, E.~Wissam, et~al.
\newblock Fine-tuning {GPT-5} for {GPU} kernel generation.
\newblock \emph{arXiv preprint arXiv:2602.11000}, 2026.

\bibitem{bai2026dice}
H.~Bai, L.~Kong, X.~Chen, et~al.
\newblock {DICE}: Diffusion large language models excel at generating {CUDA} kernels.
\newblock \emph{arXiv preprint arXiv:2602.11715}, 2026.

\bibitem{chen2026avo}
T.~Chen, Z.~Ye, B.~Xu, et~al.
\newblock {AVO}: Agentic variation operators for autonomous evolutionary search.
\newblock \emph{arXiv preprint arXiv:2603.24517}, 2026.

\bibitem{cao2026ksearch}
S.~Cao, Z.~Mao, J.~E.~Gonzalez, and I.~Stoica.
\newblock {K-Search}: {LLM} kernel generation via co-evolving intrinsic world model.
\newblock \emph{arXiv preprint arXiv:2602.19128}, 2026.

\bibitem{bhattacharjee2026optiml}
A.~Bhattacharjee, H.~Ping, S.~V.~Le, et~al.
\newblock {OptiML}: An end-to-end framework for program synthesis and {CUDA} kernel optimization.
\newblock \emph{arXiv preprint arXiv:2602.12305}, 2026.

\bibitem{dong2026kernelblaster}
K.~S.~Dong, S.~Modi, D.~Nikiforov, et~al.
\newblock {KernelBlaster}: Continual cross-task {CUDA} optimization via memory-augmented in-context reinforcement learning.
\newblock \emph{arXiv preprint arXiv:2602.14293}, 2026.

\bibitem{cao2026ascendkernelgen}
X.~Cao, J.~Zhai, P.~Li, et~al.
\newblock {AscendKernelGen}: A systematic study of {LLM}-based kernel generation for neural processing units.
\newblock \emph{arXiv preprint arXiv:2601.07160}, 2026.

\bibitem{wen2026ascendcraft}
Z.~Wen, S.~Shao, Z.~Li, et~al.
\newblock {AscendCraft}: Automatic {Ascend} {NPU} kernel generation via {DSL}-guided transcompilation.
\newblock \emph{arXiv preprint arXiv:2601.22760}, 2026.

\bibitem{verdoolaege2012equiv}
S.~Verdoolaege, G.~Janssens, and M.~Bruynooghe.
\newblock Equivalence checking of static affine programs using widening to handle recurrences.
\newblock In \emph{Proc. CAV}, 2012.

\bibitem{barthou2012widening}
D.~Barthou, P.~Music, and A.~Cohen.
\newblock Equivalence checking of static affine programs.
\newblock \emph{ACM SIGPLAN Notices}, 2002.

\bibitem{lefevre2023verified}
C.~Lef\`evre, A.~Music, and L.-N.~Pouchet.
\newblock Verified code generation for the polyhedral model.
\newblock In \emph{Proc. PLDI}, 2023.

\end{thebibliography}

\end{document}
